\SourceLeaf{453}%
\Needspace{6\baselineskip}
\[
\dfrac{3014{,}48021}{3014{,}5757}=\operatorname{num}(\log=-0{,}00001376)
\]
\noindent für die Basislänge nach \textsc{Peters} und für die von \textsc{Gauss} in der ‘Breitenbestimmung’ wie oben angegebenen Erddimensionen,
\par

\Needspace{6\baselineskip}
\[
\dfrac{3014{,}451}{3014{,}5757}=\operatorname{num}(\log=-0{,}00001797)
\]
\noindent für die Basislänge nach \textsc{Peters} und \textsc{Andrae} und für die von \textsc{Gauss} in der ‘Breitenbestimmung’ wie oben angegebenen Erddimensionen,
\par

\Needspace{6\baselineskip}
\[
\dfrac{3014{,}48021}{3014{,}5757}.\dfrac{57009{,}746}{57009{,}7584}=\operatorname{num}(\log=-0{,}00001386)
\]
\noindent für die Basislänge nach \textsc{Peters} und für \textsc{Walbeck}’s Erd dimensionen,
\par

\Needspace{6\baselineskip}
\[
\dfrac{3014{,}451}{3014{,}5757}.\dfrac{57009{,}746}{57009{,}7584}=\operatorname{num}(\log=-0{,}00001806)
\]
\noindent für die Basislänge nach \textsc{Peters} und \textsc{Andrae} und für \textsc{Walbeck}’s Erddimensionen,
\par

\Needspace{6\baselineskip}
\[
\dfrac{3014{,}48021}{3014{,}5757}.\dfrac{57009{,}746}{57008{,}551}=\operatorname{num}(\log=-0{,}00000466)
\]
\noindent für die Basislänge nach \textsc{Peters} und für \textsc{Schmidt}’s IV. Erddimensionen,
\par

\Needspace{6\baselineskip}
\[
\dfrac{3014{,}451}{3014{,}5757}.\dfrac{57009{,}746}{57008{,}551}=\operatorname{num}(\log=-0{,}00000887)
\]
\noindent für die Basislänge nach \textsc{Peters} und \textsc{Andrae} und für \textsc{Schmidt}’s IV. Erddimensionen
\par

\Needspace{6\baselineskip}
\[
\dfrac{3014{,}48021}{3014{,}5757}.\dfrac{57009{,}746}{57013{,}109}=\operatorname{num}(\log=-0{,}00003938)
\]
\noindent für die Basislänge nach \textsc{Peters} und für \textsc{Bessel}’s Erddimensionen,
\par

\Needspace{6\baselineskip}
\[
\dfrac{3014{,}451}{3014{,}5757}.\dfrac{57009{,}746}{57013{,}109}=\operatorname{num}(\log=-0{,}00004359)
\]
\noindent für die Basislänge nach \textsc{Peters} und \textsc{Andrae} und für \textsc{Bessel}’s Erddimensionen.
\par

Die von \textsc{Schmidt} und die von \textsc{Bessel} berechneten Erddimensionen setzen die Längenangabe von \textsc{Schumacher} über dessen Braacker Basis voraus, eine neue Berechnung der von ihnen in Betracht gezogenen Gradmessungen würde bei dieser berichtigten Basislänge etwas abweichende Zahlen für die Erddimensionen ergeben, die aber durch die bald zu erwartende Beendigung mehrer neuen Gradmessungen auch in kurzer Zeit durch bessere Bestimmungen ersetzt werden müssen.

Die hier im Abdruck aus den Partial-Verzeichnissen noch besonders aufgenommenen Coordinaten, sind entweder dieselben wie im General-Verzeichniss oder beruhen auf weniger genauen Bestimmungen. können aber zur Erläuterung der nachfolgenden ‘Abrisse’ dienen. In \textsc{Gauss} Nachlass befinden sich von den Partial-Verzeichnissen nur Nro. 1 bis 11. Eine neue Vergleichung ergab mir die Berichtigungen:
\begin{center}\small\begin{tabular}{@{}p{.39\linewidth}rrp{.16\linewidth}@{}}
im General-Verzeichniss steht:&\(-26619{,}9\)&\(-12689{,}5\)&Lauenberge\\
im Partial-Verzeichniss (3) steht:&\(-26619{,}9\)&\(+12689{,}5\)&Lauenberge
\end{tabular}\end{center}
\begin{center}\begin{tabular}{@{}rrp{.55\linewidth}@{}}
\(-233491{,}171\)&\(+181317{,}782\)&Norden Thürmchen auf hoher Kirche. Nr. 12.\\
\(-249451{,}172\)&\(+162280{,}007\)&Langeoog F. J. \textsc{Pauls} Haus östlicher Giebelstock. Nr. 12.
\end{tabular}\end{center}

Zur leichtern Wiedererkennung der in dem Coordinaten-Verzeichniss angegebenen Punkte kann man die auf diese Vermessung gegründete ‘\textsc{Papen}’sche Karte vom Königreich Hannover’ mit Vortheil benutzen.

Die Überschriften \(+\) südlich und \(+\) westlich habe ich, um dem Rechner ein Missverstehen der Zeichen sicherer vermeiden zu lassen, hinzugefügt.
\begin{flushright}\textsc{Schering.}\end{flushright}
\begin{center}\rule{.23\linewidth}{.3pt}\end{center}
